\subsection{Costos y beneficios}

A continuación se presentan los costos y beneficios que se tienen planeados se puedan generar por la implementación del proyecto. Todas las cantidades monetarias se presentan en pesos mexicanos (MXN).

\subsubsection{Costos}

Desarrollo del proyecto:

\medskip

\noindent
\begin{tabularx}{\textwidth}{>{\raggedright\arraybackslash}X>{\raggedright\arraybackslash}X}
  \toprule
  Concepto & Costo \\
  \midrule
  Horas del equipo y de los expertos/colaboradores & \$0 (sin remuneración) \\
  PC de entrenamiento e iPhones & \$0 (propios) \\
  Software (Python, TensorFlow, Jupyter, Git\dots) & \$0 (código abierto) \\
  Acceso al CAETEC & \$0 \\
  Gasolina de la visita extra & \$150 \\
  \midrule
  Total & \$150 \\
  \bottomrule
\end{tabularx}

\medskip

Riesgos que pueden generar costos extra:

\medskip

\noindent
\begin{tabularx}{\textwidth}{>{\raggedright\arraybackslash}X>{\raggedright\arraybackslash}X>{\raggedright\arraybackslash}X}
  \toprule
  Riesgo & Nivel de riesgo & Costo si ocurre \\
  \midrule
  Dataset chico/desbalanceado & 85\% & \$300 (2 visitas) \\
  Calidad irregular de fotos & 80\% & \$150 (1 visita) \\
  \midrule
  Total & \multicolumn{2}{r}{\$450} \\
  \bottomrule
\end{tabularx}

\medskip

Teniendo en cuenta el análisis anterior, se estima que el costo total del proyecto esté en un rango de \$600 - \$700 pesos mexicanos.

Costos de despliegue formal en campo:

Estos costos son los que se tienen previstos que habría si se decidiera implementar este proyecto a mayor escala ya en el CAETEC. Cabe mencionar que estos costos no forman parte del proyecto a realizar y solo funcionan como un mero escenario hipotético para tener un panorama de lo que cuesta aproximadamente implementar un proyecto de este tipo ya de manera formal.

\medskip

\noindent
\begin{tabularx}{\textwidth}{>{\raggedright\arraybackslash}X>{\raggedright\arraybackslash}X}
  \toprule
  Concepto & Costo \\
  \midrule
  Cámara 3D & \$10,000 \\
  Instalación y red & \$12,000 \\
  Equipo de cómputo Edge Inference & \$8,000 \\
  Mantenimiento del hardware (cada año) & \$3,000 \\
  \midrule
  Total & \$33,000 \\
  \bottomrule
\end{tabularx}

\subsubsection{Beneficios anuales si el proyecto es exitoso (establo de 160 vacas)}

Para el cálculo de los beneficios, partimos de los siguientes supuestos:

Cabe recalcar que todos estos porcentajes y montos son meras estimaciones del equipo.

\medskip

\noindent
\begin{tabularx}{\textwidth}{>{\raggedright\arraybackslash}X>{\raggedright\arraybackslash}X}
  \toprule
  Beneficio & Supuesto \\
  \midrule
  Mano de obra & Calificar el BCS a mano toma 3 min por vaca, una vez al mes, a \$120/h \\
  \addlinespace
  Salud y reproducción & 5 \% del establo tiene un problema ligado al BCS (como cetosis por ejemplo), la detección temprana evita el 30 \% y cada caso cuesta \$8,000 \\
  \addlinespace
  Concentrado & 8 kg por vaca al día a \$7/kg, con un ahorro de 1 \% \\
  \addlinespace
  Colas & Cada vaca produce 0.2 L más al día y la leche se vende a \$8.50/L \\
  \bottomrule
\end{tabularx}

\bigskip

\noindent
\begin{tabularx}{\textwidth}{>{\raggedright\arraybackslash}X>{\raggedright\arraybackslash}X>{\raggedright\arraybackslash}X}
  \toprule
  Beneficio & Cómo se calcula & Por año \\
  \midrule
  Ahorro de mano de obra & Calificar el BCS a mano toma 3 min por vaca al mes: 160 $\times$ 3 min $\times$ 12 = 96 h/año, a \$120/h & \$11,500 \\
  \addlinespace
  Menos problemas de salud y reproducción & 5 \% del establo tiene un problema ligado al BCS; detectarlo a tiempo evita el 30 \% ($\approx$ 2.4 casos) a \textasciitilde\$8,000 cada uno & \$19,200 \\
  \addlinespace
  Ajuste del concentrado & El campo gasta \textasciitilde\$3.27 M/año en concentrado (160 vacas $\times$ 8 kg $\times$ \$7 $\times$ 365); dar la ración según el BCS ahorra 1 \% & \$32,700 \\
  \midrule
  Subtotal BCS & \multicolumn{2}{r}{\$63,400} \\
  \midrule
  Solución de las colas & Si bajan las esperas, el robot se aprovecha mejor y cada vaca da \textasciitilde 0.2 L más al día: 160 $\times$ 0.2 $\times$ 365 = 11,680 L $\times$ \$8.50 & \$99,300 \\
  \bottomrule
\end{tabularx}

\medskip

El análisis costo-beneficio indica que el proyecto es económicamente viable y de muy bajo riesgo financiero. Nuestro costo real de desarrollo es mínimo, solo existen gastos de transporte y gasolina si los riesgos del dataset obligaran a hacer visitas adicionales. Esto se debe a que el tiempo del equipo y de los expertos del CAETEC es una aportación sin remuneración, y el equipo, el software y el acceso al campo no tienen costo.

La mayor parte del beneficio proviene del ajuste de concentrado según el BCS y de la posible solución al problema de las colas, aunque esta última es la parte más incierta por el riesgo de falta de tiempo.
